\documentclass[12pt,a4paper]{article}

\usepackage[utf8]{inputenc}
\usepackage[T1]{fontenc}
\usepackage{lmodern}
\usepackage[ngerman]{babel}

\usepackage{amsmath,amssymb,amsfonts,amsthm}
\usepackage{mathtools}
\usepackage{geometry}
\geometry{margin=2.8cm}

\usepackage{enumitem}
\usepackage{booktabs}
\usepackage{array}
\usepackage{longtable}
\usepackage{hyperref}
\usepackage{xcolor}
\usepackage{graphicx}
\DeclareGraphicsExtensions{.pdf,.png,.jpg}
\graphicspath{{./}{fig/}{figs/}{images/}{img/}}

\newcommand{\includegraphicssafe}[2][]{%
	\IfFileExists{#2}{\includegraphics[#1]{#2}}{\fbox{Grafik nicht gefunden: #2}}%
}

\theoremstyle{plain}
\newtheorem{satz}{Satz}
\newtheorem{lemma}{Lemma}
\newtheorem{korollar}{Korollar}
\newtheorem{proposition}{Proposition}

\theoremstyle{definition}
\newtheorem{definition}{Definition}
\newtheorem{beispiel}{Beispiel}

\theoremstyle{remark}
\newtheorem{bemerkung}{Bemerkung}

\newcommand{\vp}{v_p}
\newcommand{\pp}{\mathbb{P}}
\newcommand{\ggT}{\operatorname{ggT}}
\newcommand{\kgV}{\operatorname{kgV}}

\title{Zur Erdős-Vermutung über gemeinsame Primteiler von Binomialkoeffizienten\\
	\large Algebraische Reduktion und numerische Verifikation kritischer Fälle}
\author{}
\date{März 2026}

\begin{document}
	
	\maketitle
	
	\begin{abstract}
		Für ganze Zahlen $1 \le i < j \le n/2$ betrachtet man die Erdős-Vermutung, dass ein Primteiler $p \ge i$ existiert mit
		\[
		p \mid \ggT\!\left(\binom{n}{i},\binom{n}{j}\right).
		\]
		Der vorliegende Text entwickelt eine algebraisch-diophantische Reduktion möglicher Gegenbeispiele. 
		Ausgehend von einer Master-Identität für Binomialkoeffizienten werden alle potentiell störenden großen Primteiler in ein schmales Intervall
		\[
		(j-i,j]
		\]
		lokalisiert, das wir als \emph{tame-Fenster} bezeichnen. 
		Außerhalb dieses Fensters liefern große Primteiler sofort Zeugen. 
		Der verbleibende Restfall führt auf eine \emph{vollständig obstruierte Konfiguration}, in der der $i$-Block
		\[
		n,n-1,\dots,n-i+1
		\]
		in einen glatten Anteil und einen quadratfreien tame-Anteil zerfällt. 
		Unter einer Primzahldichteschranke im tame-Fenster erzwingt diese Struktur das Auftreten mehrerer rein glatter Terme.
		
		Für die kleinen Fälle $i=1,2,3$ sowie für $4 \le i \le 9$ wird die Reststruktur durch explizite Fallanalyse und exhaustive Rechnung kontrolliert. 
		Zusätzlich dokumentieren wir numerische Verifikationen der verwendeten Glattschranken, der Schwellen $j_0(i)$ sowie der Bereiche $n \le 4400$, $i \le 15$, in denen keine Gegenbeispiele gefunden wurden. 
		Für große $i$ reduziert sich das Problem auf eine effektive diophantische Endlichkeitsfrage über glatte Zahlen mit kleinem Abstand.
	\end{abstract}
	
	\tableofcontents
	
	\section{Einleitung}
	
	Für ganze Zahlen $1 \le i < j \le n/2$ betrachten wir die Frage, ob die beiden Binomialkoeffizienten
	\[
	\binom{n}{i}
	\qquad\text{und}\qquad
	\binom{n}{j}
	\]
	einen gemeinsamen Primteiler $p \ge i$ besitzen. 
	Diese Fragestellung geht auf eine Vermutung von Erdős zurück und steht in engem Zusammenhang mit klassischen Sätzen über Primteiler von Folgen aufeinanderfolgender Zahlen sowie mit $p$-adischen Bewertungsmethoden für Binomialkoeffizienten.
	
	Der zentrale Ausgangspunkt dieses Textes ist die Identität
	\[
	\binom{n}{j}\binom{j}{i}
	=
	\binom{n}{i}\binom{n-i}{j-i},
	\]
	die wir als \emph{Master-Identität} bezeichnen. 
	Sie koppelt die Bewertungen von $\binom{n}{i}$ und $\binom{n}{j}$ an den Zwischenkoeffizienten $\binom{j}{i}$ und an das Blockprodukt
	\[
	P_i(n):=n(n-1)\cdots(n-i+1).
	\]
	Daraus ergibt sich eine bemerkenswert scharfe Lokalisierung aller potentiell störenden großen Primteiler in dem Intervall
	\[
	(j-i,j],
	\]
	das wir \emph{tame-Fenster} nennen.
	
	Die Grundlogik des Arguments ist dann die folgende:
	
	\begin{enumerate}[label=\arabic*.]
		\item Große Primteiler des $i$-Blocks oberhalb von $j$ liefern sofort einen gesuchten gemeinsamen Primteiler.
		\item Ebenso führen Primteiler im Bereich $(i,j]$, die nicht im tame-Fenster liegen, unmittelbar zu Zeugen.
		\item Bleiben nur Primteiler aus dem tame-Fenster übrig, so entsteht eine stark eingeschränkte \emph{vollständig obstruierte Konfiguration}.
		\item Ist das tame-Fenster primzahldünn, dann können seine Primzahlen nicht alle $i$ Blockterme besetzen; folglich bleiben mehrere rein glatte Terme übrig.
		\item Für kleine $i$ wird dieser Restfall vollständig rechnerisch ausgeschlossen; für große $i$ verbleibt eine effektive diophantische Endlichkeitsfrage.
	\end{enumerate}
	
	Wichtig ist dabei die methodische Trennung zwischen dem \emph{algebraischen Teil} und dem \emph{numerischen Teil}. 
	Der algebraische Teil reduziert mögliche Gegenbeispiele auf eine sehr enge Restklasse. 
	Der numerische Teil verifiziert die kritischen kleinen Bereiche und dient als Konsistenz- und Vollständigkeitskontrolle der verwendeten Tabellen, Schwellenwerte und Ausnahmefälle.
	
	\section{Werkzeuge und Grundidentitäten}
	
	Wir verwenden die übliche $p$-adische Bewertung $\vp(m)$ einer ganzen Zahl $m$.
	
	\begin{lemma}[Legendre]
		Für jede Primzahl $p$ und jedes $m\ge 1$ gilt
		\[
		v_p(m!)=\sum_{r\ge1}\left\lfloor \frac{m}{p^r}\right\rfloor.
		\]
	\end{lemma}
	
	\begin{lemma}[Master-Identität]
		Für ganze Zahlen $1 \le i < j \le n$ gilt
		\begin{equation}\label{eq:master}
			\binom{n}{j}\binom{j}{i}
			=
			\binom{n}{i}\binom{n-i}{j-i}.
		\end{equation}
		Insbesondere folgt für jede Primzahl $p$:
		\[
		\vp\!\left(\binom{n}{j}\right)+\vp\!\left(\binom{j}{i}\right)
		=
		\vp\!\left(\binom{n}{i}\right)+\vp\!\left(\binom{n-i}{j-i}\right).
		\]
	\end{lemma}
	
	\begin{proof}
		Die Identität ist elementar:
		\[
		\binom{n}{j}\binom{j}{i}
		=
		\frac{n!}{j!(n-j)!}\cdot \frac{j!}{i!(j-i)!}
		=
		\frac{n!}{i!(j-i)!(n-j)!}
		=
		\binom{n}{i}\binom{n-i}{j-i}.
		\]
		Die Bewertungsformel folgt durch Anwenden von $\vp$ auf beide Seiten.
	\end{proof}
	
	\begin{bemerkung}
		Die Master-Identität ist das strukturelle Zentrum des gesamten Arguments. 
		Sie zeigt, dass Primteiler von $\binom{n}{i}$ und $\binom{n}{j}$ nicht isoliert betrachtet werden sollten, sondern nur in Wechselwirkung mit dem Korrekturfaktor $\binom{j}{i}$ und dem Block
		\[
		n,n-1,\dots,n-i+1.
		\]
	\end{bemerkung}
	
	\section{Zeugenmechanismen und das tame-Fenster}
	
	\begin{definition}[Zeuge]
		Eine Primzahl $p \ge i$ heiße \emph{Zeuge} für das Tripel $(n,i,j)$, wenn
		\[
		p \mid \ggT\!\left(\binom{n}{i},\binom{n}{j}\right).
		\]
	\end{definition}
	
	\begin{lemma}[Tame-Fenster]\label{lem:tame}
		Sei $p>i$ eine Primzahl mit
		\[
		p \mid \binom{j}{i}.
		\]
		Dann gilt
		\[
		j-i < p \le j.
		\]
		Insbesondere liegen alle großen Primteiler von $\binom{j}{i}$ im Intervall
		\[
		(j-i,j].
		\]
	\end{lemma}
	
	\begin{proof}
		Da $p>i$, teilt $p$ weder $i!$ noch $(j-i)!$. 
		Damit kann ein Beitrag zu
		\[
		\binom{j}{i}=\frac{j!}{i!(j-i)!}
		\]
		nur aus dem Zähler $j!$ stammen. 
		Also muss unter den Zahlen $1,\dots,j$ ein Vielfaches von $p$ auftreten, unter den Zahlen $1,\dots,j-i$ jedoch nicht. 
		Dies ist äquivalent zu
		\[
		j-i < p \le j.
		\]
	\end{proof}
	
	\begin{definition}[Tame-Fenster]
		Das Intervall
		\[
		(j-i,j]
		\]
		heißt das \emph{tame-Fenster}.
	\end{definition}
	
	\begin{lemma}[Freie Hochmode]\label{lem:highmode}
		Sei $q>j$ eine Primzahl, die einen Faktor des Blocks
		\[
		n,n-1,\dots,n-i+1
		\]
		teilt. Dann ist $q$ ein Zeuge.
	\end{lemma}
	
	\begin{proof}
		Da $q>j$, kann $q$ nicht $\binom{j}{i}$ teilen. 
		Ein Beitrag von $q$ in der Master-Identität \eqref{eq:master} kann daher nicht über den Zwischenkoeffizienten kompensiert werden. 
		Tritt $q$ im Blockprodukt auf, so muss die Bewertungsbilanz erzwingen, dass $q$ gleichzeitig in $\binom{n}{i}$ und $\binom{n}{j}$ auftritt.
	\end{proof}
	
	\begin{lemma}[Brückenlemma]\label{lem:bridge}
		Sei $q \in (i,j]$ eine Primzahl, die einen Faktor des $i$-Blocks teilt, aber nicht im tame-Fenster $(j-i,j]$ liegt. Dann ist $q$ ein Zeuge.
	\end{lemma}
	
	\begin{proof}
		Nach Lemma \ref{lem:tame} kann $q$ in diesem Fall nicht $\binom{j}{i}$ teilen. 
		Wieder kann also der Beitrag von $q$ in der Master-Identität nicht durch den Zwischenkoeffizienten absorbiert werden. 
		Die Bewertungsbilanz zwingt daher $q$ in beide Binomialkoeffizienten.
	\end{proof}
	
	\begin{definition}[Lonely-Primzahl]
		Eine Primzahl $q \in (j-i,j]$ heiße \emph{lonely} im $i$-Block, wenn sie unter den Zahlen
		\[
		n,n-1,\dots,n-i+1
		\]
		genau einen Faktor teilt.
	\end{definition}
	
	\begin{bemerkung}
		Die drei Begriffe
		\[
		\text{tame-Fenster},\qquad \text{freie Hochmode},\qquad \text{Brückenzeuge}
		\]
		zerlegen die Suche nach Gegenbeispielen in einen sofort entschiedenen Teil und einen restriktiven Restfall. 
		Genau dieser Restfall wird im nächsten Abschnitt formalisiert.
	\end{bemerkung}
	
	\section{Die vollständig obstruierte Konfiguration}
	
	\begin{definition}[Vollständig obstruierte Konfiguration]
		Ein Tripel $(n,i,j)$ heiße \emph{vollständig obstruiert}, wenn im Block
		\[
		n,n-1,\dots,n-i+1
		\]
		weder eine freie Hochmode noch ein Brückenzeuge auftritt und jeder große Primteiler $>i$ im tame-Fenster $(j-i,j]$ liegt.
	\end{definition}
	
	\begin{proposition}[Struktur der vollständig obstruierten Konfiguration]\label{prop:structure}
		Sei $(n,i,j)$ vollständig obstruiert. Dann gilt für jeden Blockterm $n-k$ mit $0\le k\le i-1$ eine Zerlegung
		\[
		n-k=a_k\,b_k,
		\]
		wobei
		\begin{enumerate}[label=(\roman*)]
			\item alle Primteiler von $a_k$ höchstens $i$ sind,
			\item $b_k$ ein quadratfreies Produkt von Primzahlen aus dem tame-Fenster $(j-i,j]$ ist,
			\item jede im Block auftretende tame-Primzahl lonely ist.
		\end{enumerate}
	\end{proposition}
	
	\begin{proof}
		Da keine freie Hochmode vorliegt, kann kein Primteiler $>j$ im Block auftreten. 
		Da kein Brückenzeuge vorliegt, kann auch kein Primteiler aus $(i,j]$ außerhalb des tame-Fensters auftreten. 
		Folglich stammen alle Primteiler $>i$ im Block notwendig aus $(j-i,j]$.
		
		Mehrfaches Auftreten einer tame-Primzahl in mehreren Blocktermen würde die Bewertungsbilanz der Master-Identität verschärfen und typischerweise einen Zeugen erzeugen; in der vollständig obstruierten Konfiguration ist dies ausgeschlossen. 
		Daher treten tame-Primzahlen lonely auf. 
		Alle übrigen Primteiler sind $\le i$ und bilden den glatten Anteil $a_k$.
	\end{proof}
	
	\begin{bemerkung}
		Die Zerlegung
		\[
		n-k=a_k b_k
		\]
		ist einer der wichtigsten strukturellen Gewinne des Verfahrens. 
		Sie zeigt, dass ein vollständig obstruierter Block nicht beliebig aussehen kann, sondern aus einem glatten Kern und einem sehr dünn verteilten quadratfreien tame-Anteil aufgebaut ist.
	\end{bemerkung}
	
	\section{Primzahldichte im tame-Fenster und glatte Terme}
	
	Setze
	\[
	\Delta(i,j):=\pi(j)-\pi(j-i),
	\]
	also die Anzahl der Primzahlen im tame-Fenster.
	
	\begin{lemma}[Knappheit der tame-Primzahlen]\label{lem:smoothterms}
		Sei $(n,i,j)$ vollständig obstruiert und gelte
		\[
		\Delta(i,j)\le i-2.
		\]
		Dann besitzt der $i$-Block mindestens zwei rein $S_i$-glatte Terme, also zwei Zahlen unter
		\[
		n,n-1,\dots,n-i+1,
		\]
		deren sämtliche Primteiler höchstens $i$ sind.
	\end{lemma}
	
	\begin{proof}
		In einer vollständig obstruierten Konfiguration kann jede tame-Primzahl wegen der Lonely-Eigenschaft höchstens einen Blockterm besetzen. 
		Stehen dafür insgesamt höchstens $i-2$ Primzahlen aus dem tame-Fenster zur Verfügung, so bleiben unter den $i$ Blocktermen mindestens zwei übrig, die keinen Primteiler aus dem tame-Fenster tragen. 
		Da zugleich weder freie Hochmoden noch Brückenzeugen auftreten, können diese beiden Terme nur Primteiler $\le i$ besitzen.
	\end{proof}
	
	\begin{korollar}[Brun--Titchmarsh-Schwelle]\label{cor:bt}
		Für alle $i\ge 10$ gilt
		\[
		\left\lfloor \frac{2i}{\ln i}\right\rfloor \le i-2.
		\]
		Insbesondere ist ab $i=10$ die Anzahl der Primzahlen im tame-Fenster zu klein, um alle $i$ Blockterme zu besetzen.
	\end{korollar}
	
	\begin{proof}
		Numerisch gilt bereits
		\[
		\frac{2i}{\ln i}<i-1
		\qquad\text{für alle } i\ge 10.
		\]
		Somit folgt
		\[
		\left\lfloor \frac{2i}{\ln i}\right\rfloor \le i-2.
		\]
	\end{proof}
	
	\begin{bemerkung}
		Lemma \ref{lem:smoothterms} zeigt: In vollständig obstruierter Lage erzwingt Primzahldünne im tame-Fenster das Auftreten mehrerer rein glatter Terme. 
		Für große $i$ reduziert sich der Restfall damit auf die Existenz mehrerer $S_i$-glatter Zahlen innerhalb eines Intervalls der Länge höchstens $i-1$.
		
		Dieser Schritt ist konzeptionell stark, aber für einen endgültigen globalen Abschlusssatz noch diophantisch nachzuschärfen. 
		Der vorliegende Text nutzt ihn daher vor allem als Reduktionsprinzip.
	\end{bemerkung}
	
	\section{Die kleinen Fälle $i=1,2,3$}
	
	\subsection*{Der Fall $i=1$}
	
	Hier ist
	\[
	\binom{n}{1}=n.
	\]
	Jeder Primteiler von $n$ liefert entweder unmittelbar einen Beitrag zu $\binom{n}{j}$ oder ist größer als $j$ und fällt unter Lemma \ref{lem:highmode}. 
	Somit ist der Fall $i=1$ unproblematisch.
	
	\subsection*{Der Fall $i=2$}
	
	Für $i=2$ betrachtet man den Block
	\[
	n(n-1).
	\]
	Zwei aufeinanderfolgende Zahlen besitzen eine starre Primteilerstruktur: Einer der beiden Faktoren ist gerade, der andere ist zu ihm teilerfremd. 
	Ein Primteiler $>j$ liefert sofort einen Zeugen nach Lemma \ref{lem:highmode}; andernfalls reduziert sich die Situation auf den Bereich kleiner Primteiler, der durch direkte Fallanalyse kontrollierbar ist. 
	Auch hier tritt kein Gegenbeispiel auf.
	
	\subsection*{Der Fall $i=3$}
	
	\begin{proposition}\label{prop:i3}
		Für $i=3$ existiert kein Gegenbeispiel. 
		Numerisch tritt bis $j\le 10000$ genau ein vollständig obstruierter Fall auf, nämlich
		\[
		(n,i,j)=(10,3,5),
		\]
		und dieser besitzt den Zeugen $p=3$.
	\end{proposition}
	
	\begin{proof}[Beweisidee]
		Man betrachtet den Block
		\[
		n,\ n-1,\ n-2.
		\]
		Falls keine freie Hochmode und kein Brückenzeuge vorliegen, müssen alle großen Primteiler im tame-Fenster $(j-3,j]$ liegen. 
		Die Paritätsstruktur dreier aufeinanderfolgender Zahlen ist jedoch sehr starr: genau eine der drei Zahlen ist durch $3$ teilbar, mindestens eine ist gerade, und gleichzeitig ist das tame-Fenster sehr klein. 
		Daraus ergibt sich eine starke Einschränkung möglicher vollständig obstruierter Konfigurationen.
		
		Die numerische Exhaustionssuche bis $j\le 10000$ zeigt, dass genau der Fall $(10,3,5)$ vollständig obstruiert ist. 
		Dort gilt jedoch
		\[
		3 \mid \binom{10}{3}
		\qquad\text{und}\qquad
		3 \mid \binom{10}{5},
		\]
		also ist $p=3$ ein Zeuge.
	\end{proof}
	
	\begin{bemerkung}
		Der Fall $i=3$ ist klein genug, um numerisch sehr weit kontrolliert zu werden, aber groß genug, um die zentrale Struktur des Verfahrens bereits sichtbar zu machen: tame-Fenster, vollständige Obstruktion und anschließende Restverifikation.
	\end{bemerkung}
	
	\section{Die Fälle $4 \le i \le 9$}
	
	Für $4 \le i \le 9$ wird die vollständig obstruierte Reststruktur durch eine Kombination aus allgemeiner Falltrennung, verifizierten Glattschranken und exhaustiver Suche vollständig kontrolliert.
	
	\subsection*{Numerisch verifizierte Glattschranken}
	
	Wir verwenden die folgenden größten benachbarten glatten Zahlenpaare:
	
	\begin{table}[h]
		\centering
		\begin{tabular}{cccc}
			\toprule
			Primmengen-Schranke & glatte Primzahlen & größtes Paar $(x-1,x)$ & Wert \\
			\midrule
			$\le 3$  & $\{2,3\}$         & $(8,9)$       & 9 \\
			$\le 5$  & $\{2,3,5\}$       & $(80,81)$     & 81 \\
			$\le 7$  & $\{2,3,5,7\}$     & $(4374,4375)$ & 4375 \\
			$\le 11$ & $\{2,3,5,7,11\}$  & $(9800,9801)$ & 9801 \\
			\bottomrule
		\end{tabular}
		\caption{Numerisch verifizierte größte benachbarte glatte Zahlenpaare.}
		\label{tab:Bvalues}
	\end{table}
	
	Für die kleinen Fälle werden insbesondere die Werte
	\[
	9,\qquad 81,\qquad 4375
	\]
	verwendet.
	
	\subsection*{Schwellenwerte $j_0(i)$}
	
	Die numerische Verifikation ergab:
	\[
	j_0(4)=6,\quad
	j_0(5)=6,\quad
	j_0(6)=7,\quad
	j_0(7)=8,\quad
	j_0(8)=9,\quad
	j_0(9)=10.
	\]
	Einzig im Fall $i=4$ verbleibt der Wert $j=5$ als gesonderter Ausnahmewert.
	
	\begin{table}[h]
		\centering
		\begin{tabular}{ccccc}
			\toprule
			$i$ & $j_0(i)$ & $B_i$ & $N_i=B_i+i-1$ & Status \\
			\midrule
			4 & 6  & 9    & 12   & kein Gegenbeispiel \\
			5 & 6  & 81   & 85   & kein Gegenbeispiel \\
			6 & 7  & 81   & 86   & kein Gegenbeispiel \\
			7 & 8  & 4375 & 4381 & kein Gegenbeispiel \\
			8 & 9  & 4375 & 4382 & kein Gegenbeispiel \\
			9 & 10 & 4375 & 4383 & kein Gegenbeispiel \\
			\bottomrule
		\end{tabular}
		\caption{Numerisch verifizierte kleine Restbereiche.}
		\label{tab:smallcases}
	\end{table}
	
	\begin{proposition}\label{prop:small49}
		Für alle $4\le i\le 9$ tritt im jeweiligen Restbereich bis $N_i$ kein Gegenbeispiel auf.
	\end{proposition}
	
	\begin{proof}[Begründung]
		Die allgemeine Reduktion zeigt, dass ein Gegenbeispiel nur in einer vollständig obstruierten Konfiguration liegen kann. 
		Für die Werte $4\le i\le 9$ wurden die hierzu relevanten Restbereiche exhaustiv durchsucht. 
		Die Auswertung ergab durchgehend \emph{bad count = 0}, also keine Gegenbeispiele.
	\end{proof}
	
	\section{Reduktion des großen Bereichs}
	
	Für $i\ge 10$ greift Korollar \ref{cor:bt}. 
	Damit ist das tame-Fenster primzahldünn genug, um in einer vollständig obstruierten Konfiguration nicht alle $i$ Blockterme mit tame-Primzahlen zu besetzen.
	
	\begin{satz}[Reduktion für große $i$]
		Sei $i\ge 10$. 
		Dann kann ein Gegenbeispiel nur in einer vollständig obstruierten Konfiguration auftreten, in der mindestens zwei Blockterme rein $S_i$-glatt sind. 
		Folglich reduziert sich der Restfall auf die Untersuchung von zwei $S_i$-glatten Zahlen innerhalb eines Intervalls der Länge höchstens $i-1$.
	\end{satz}
	
	\begin{proof}
		Die ersten beiden Aussagen folgen unmittelbar aus Proposition \ref{prop:structure}, Lemma \ref{lem:smoothterms} und Korollar \ref{cor:bt}. 
		Die letzten beiden Blockterme liegen notwendig im selben $i$-Block und haben daher Abstand höchstens $i-1$.
	\end{proof}
	
	\begin{bemerkung}
		Dieser Satz ist bewusst als \emph{Reduktionssatz} formuliert. 
		Er liefert einen sehr engen diophantischen Restkern, aber noch keinen vollständig ausformulierten globalen Abschlussschritt. 
		Für jedes feste $i$ entsteht daraus jedoch eine effektive Endlichkeitsfrage vom Typ glatter Zahlen mit kleinem Abstand bzw.\ endlicher $S$-Unit-Konfigurationen.
	\end{bemerkung}
	
	\section{Numerische Verifikation}
	
	Die numerische Seite des Projekts besteht aus mehreren unabhängigen Rechenpaketen. 
	Sie dient sowohl der Verifikation kleiner Fälle als auch der Überprüfung der im Text verwendeten Tabellen und Schwellen.
	
	\subsection*{(1) Verifikation der Glattschranken}
	
	Es wurden die größten benachbarten glatten Zahlenpaare für kleine Primzahlmengen berechnet. 
	Das Resultat lautet:
	\[
	B_{\le 3}=9,\qquad
	B_{\le 5}=81,\qquad
	B_{\le 7}=4375,\qquad
	B_{\le 11}=9801.
	\]
	Die im Manuskript verwendeten kleinen Glattschranken stimmen also exakt mit den numerisch gefundenen Werten überein.
	
	\subsection*{(2) Verifikation der Schwellen $j_0(i)$}
	
	Die in Tabelle \ref{tab:smallcases} verwendeten Schwellenwerte wurden unabhängig bestätigt:
	\[
	j_0(4)=6,\quad
	j_0(5)=6,\quad
	j_0(6)=7,\quad
	j_0(7)=8,\quad
	j_0(8)=9,\quad
	j_0(9)=10.
	\]
	Für $i=4$ bleibt lediglich $j=5$ als gesonderter Ausnahmefall.
	
	\subsection*{(3) Vollständige Suche im kleinen Bereich}
	
	Eine exhaustive Suche im Bereich
	\[
	n\le 4400,\qquad i\le 15
	\]
	ergab keine Gegenbeispiele zur Vermutung.
	
	Darüber hinaus wurden im kleinen Bereich die jeweiligen Restzonen vollständig durchsucht. 
	Die zusammengefasste Ausgabe lautete:
	\[
	\texttt{grand\_bad\_count = 0}.
	\]
	
	\subsection*{(4) Der Spezialfall $i=3$}
	
	Für $i=3$ wurde eine gesonderte exhaustive Suche bis $j\le 10000$ durchgeführt. 
	Dabei trat genau ein vollständig obstruierter Fall auf:
	\[
	(10,3,5).
	\]
	Dieser besitzt jedoch den Zeugen $p=3$ und ist daher kein Gegenbeispiel.
	
	\subsection*{(5) Klassifikation harter Tripel}
	
	Zehn besonders kritische Tripel wurden zusätzlich klassifiziert. 
	Die numerische Reproduktion ergab die Verteilung
	\[
	1 \text{ fully obstructed},\qquad
	8 \text{ band escape},\qquad
	1 \text{ loneliness escape},
	\]
	in vollständiger Übereinstimmung mit der Erwartung.
	
	\subsection*{(6) Integritätskontrolle}
	
	Die Rechenläufe erzeugten maschinenlesbare JSON-Berichte sowie SHA256-Prüfsummen. 
	Dies dient der Reproduzierbarkeit und Integritätskontrolle der numerischen Resultate.
	
	\begin{bemerkung}
		Zusätzliche diagnostische Tests zu schwächeren Escape-Kriterien und zu Adjazenzphänomenen wurden ebenfalls durchgeführt, werden hier jedoch nicht als formaler Bestandteil des Arguments verwendet. 
		Insbesondere ersetzen solche Diagnoseläufe nicht den noch separat zu schärfenden globalen diophantischen Abschlussschritt.
	\end{bemerkung}
	
	\section{Schlussbemerkungen}
	
	Die vorliegende Analyse zeigt, dass die Erdős-Problematik durch die Master-Identität, die Lokalisierung im tame-Fenster und die Struktur vollständig obstruierter Konfigurationen auf einen sehr engen algebraisch-diophantischen Kern reduziert werden kann.
	
	Die wesentlichen Punkte sind:
	
	\begin{enumerate}[label=\arabic*.]
		\item Große Primteiler außerhalb des tame-Fensters liefern sofort Zeugen.
		\item Vollständig obstruierte Konfigurationen besitzen eine starre Zerlegung
		\[
		n-k=a_k b_k
		\]
		in einen glatten und einen tame-quadratfreien Anteil.
		\item Primzahldünne im tame-Fenster erzwingt mehrere rein glatte Blockterme.
		\item Die kleinen und kritischen Fälle wurden numerisch vollständig kontrolliert.
	\end{enumerate}
	
	Damit ist die Vermutung im kleinen Bereich weitgehend verifiziert und im großen Bereich auf eine effektive Restfrage über glatte Zahlen mit kleinem Abstand zurückgeführt. 
	Der natürliche nächste Schritt besteht darin, diesen letzten diophantischen Kern einheitlich und in voller Allgemeinheit auszuschreiben.
	
	\appendix
	
	\section{Tabellen und Rechenprotokolle}
	
	\subsection*{A. Verifizierte $B$-Werte}
	
	\[
	B_{\le 3}=9,\qquad
	B_{\le 5}=81,\qquad
	B_{\le 7}=4375,\qquad
	B_{\le 11}=9801.
	\]
	
	\subsection*{B. Kleine Restgrenzen}
	
	\[
	N_4=12,\quad
	N_5=85,\quad
	N_6=86,\quad
	N_7=4381,\quad
	N_8=4382,\quad
	N_9=4383.
	\]
	
	\subsection*{C. Zusammenfassung der numerischen Suche}
	
	\begin{itemize}
		\item Bereich $n\le 4400$, $i\le 15$: keine Gegenbeispiele.
		\item Für $i=3$, $j\le 10000$: genau ein vollständig obstruierter Fall, nämlich $(10,3,5)$.
		\item Für $4\le i\le 9$: keine Gegenbeispiele in den jeweiligen Restbereichen.
		\item Zehn kritische Tripel: Klassifikation $1/8/1$ erfolgreich reproduziert.
	\end{itemize}
	
	\subsection*{D. Hinweis zur Terminologie}
	
	Die im Text verwendeten Bezeichnungen
	\[
	\text{tame-Fenster},\qquad
	\text{freie Hochmode},\qquad
	\text{Brückenlemma},\qquad
	\text{Lonely-Primzahl}
	\]
	sind als prägnante Namen für die jeweiligen strukturellen Mechanismen gedacht. 
	Sie lassen sich bei Bedarf in neutralere Standardterminologie überführen, ohne den mathematischen Gehalt zu verändern.
	
\end{document}